% Blueprint: Rouche theorem via zero counting
% Created by Numina

\begin{document}

\section{Rouché's Theorem via Zero Counting}

We work over $\mathbb{C}$. Fix $R > 0$ and write
$D = \overline{B}(0,R)$ for the closed disk, $B = B(0,R)$ for the open disk and
$S = \{z : \|z\| = R\}$ for the boundary circle. Let $f$ be meromorphic in
normal form on all of $\mathbb{C}$, let $g$ be analytic on $\mathbb{C}$, and
assume $\|g(z)\| < \|f(z)\|$ for every $z \in S$. The goal is to show that $f$
and $f+g$ have the same number of zeros in $D$, each counted with multiplicity,
where the multiplicity-counted zero count of a meromorphic $h$ is
$\sum_z \big(\operatorname{divisor}(h, D)\big)^{+}(z)$, the total mass of the
positive part of its divisor.

The strategy is the argument principle: the signed mass of the divisor of $h$ in
$D$ equals the normalized contour integral $\frac{1}{2\pi i}\oint_{C(0,R)}
\operatorname{logDeriv} h$. Applying this to $f$ and to $f+g$ and showing the two
contour integrals agree (because $(f+g)/f$ winds zero times around the origin),
together with the fact that $f$ and $f+g$ have identical pole divisors, yields
the result.

\subsection{Behaviour on the boundary circle}

\begin{lemma}[$f$ does not vanish on the circle]
    \label{lem:f-ne-zero-sphere}
    \lean{LeanEval.ComplexAnalysis.f_ne_zero_sphere}
    \leanfile{LeanEval/ComplexAnalysis/Rouche/Support.lean}
    \leanok
    For every $z \in S$ we have $f(z) \neq 0$.
\end{lemma}
\begin{proof}
    The hypothesis gives $\|f(z)\| > \|g(z)\| \ge 0$, so $\|f(z)\| > 0$, hence
    $f(z) \neq 0$.
\end{proof}

\begin{lemma}[$f+g$ does not vanish on the circle]
    \label{lem:fg-ne-zero-sphere}
    \lean{LeanEval.ComplexAnalysis.fg_ne_zero_sphere}
    \leanfile{LeanEval/ComplexAnalysis/Rouche/Support.lean}
    \leanok
    For every $z \in S$ we have $(f+g)(z) \neq 0$.
\end{lemma}
\begin{proof}
    By the reverse triangle inequality
    $\|(f+g)(z)\| \ge \|f(z)\| - \|g(z)\| > 0$ using the hypothesis
    $\|g(z)\| < \|f(z)\|$. Hence $(f+g)(z) \neq 0$.
\end{proof}

\begin{lemma}[$f$ is analytic and nonzero near the circle]
    \label{lem:f-analytic-sphere}
    \lean{LeanEval.ComplexAnalysis.f_analytic_sphere}
    \leanfile{LeanEval/ComplexAnalysis/Rouche/Support.lean}
    \uses{lem:f-ne-zero-sphere}
    \leanok
    For every $z \in S$, the meromorphic order of $f$ at $z$ is $0$; in
    particular $f$ is analytic at $z$ with $f(z) \neq 0$.
\end{lemma}
\begin{proof}
    Since $f$ is in meromorphic normal form, the order of $f$ at $z$ is zero if
    and only if $f(z) \neq 0$ (a function in normal form with a zero or pole
    takes the value $0$ there). By Lemma~\ref{lem:f-ne-zero-sphere} the value is
    nonzero, so the order is $0$, which is equivalent to $f$ being analytic at
    $z$ with nonzero value.
\end{proof}

\begin{lemma}[$f+g$ is analytic and nonzero near the circle]
    \label{lem:fg-analytic-sphere}
    \lean{LeanEval.ComplexAnalysis.fg_analytic_sphere}
    \leanfile{LeanEval/ComplexAnalysis/Rouche/Support.lean}
    \uses{lem:f-analytic-sphere, lem:fg-ne-zero-sphere}
    \leanok
    For every $z \in S$, $f+g$ is analytic at $z$ and $(f+g)(z) \neq 0$;
    consequently its meromorphic order at $z$ is $0$.
\end{lemma}
\begin{proof}
    By Lemma~\ref{lem:f-analytic-sphere} $f$ is analytic at $z$, and $g$ is
    analytic everywhere, so $f+g$ is analytic at $z$. By
    Lemma~\ref{lem:fg-ne-zero-sphere} its value is nonzero, so an analytic
    function with nonzero value has order $0$ there.
\end{proof}

\begin{lemma}[A point of order zero is off the divisor]
    \label{lem:divisor-eq-zero-of-orderZero}
    \lean{LeanEval.ComplexAnalysis.divisor_eq_zero_of_orderZero}
    \leanfile{LeanEval/ComplexAnalysis/Rouche/Support.lean}
    \leanok
    Let $h$ be meromorphic on $\mathbb{C}$ and $z \in D$ with meromorphic order
    of $h$ at $z$ equal to $0$. Then $\operatorname{divisor}(h,D)(z) = 0$.
\end{lemma}
\begin{proof}
    The divisor at a point of $D$ is the untopped meromorphic order of $h$ at
    that point. Since the order at $z$ is the integer $0$, its untopping is $0$,
    so $\operatorname{divisor}(h,D)(z) = 0$.
\end{proof}

\begin{lemma}[$f$'s divisor vanishes on the circle]
    \label{lem:divisor-f-zero-sphere}
    \lean{LeanEval.ComplexAnalysis.divisor_f_zero_sphere}
    \leanfile{LeanEval/ComplexAnalysis/Rouche/Support.lean}
    \uses{lem:f-analytic-sphere, lem:divisor-eq-zero-of-orderZero}
    \leanok
    For every $z \in S$ we have $\operatorname{divisor}(f,D)(z) = 0$.
\end{lemma}
\begin{proof}
    \uses{lem:f-analytic-sphere, lem:divisor-eq-zero-of-orderZero}
    Fix $z \in S \subseteq D$. By Lemma~\ref{lem:f-analytic-sphere} the order of
    $f$ at $z$ is $0$, so Lemma~\ref{lem:divisor-eq-zero-of-orderZero} gives
    $\operatorname{divisor}(f,D)(z) = 0$.
\end{proof}

\begin{lemma}[$f+g$'s divisor vanishes on the circle]
    \label{lem:divisor-fg-zero-sphere}
    \lean{LeanEval.ComplexAnalysis.divisor_fg_zero_sphere}
    \leanfile{LeanEval/ComplexAnalysis/Rouche/Support.lean}
    \uses{lem:fg-analytic-sphere, lem:divisor-eq-zero-of-orderZero}
    \leanok
    For every $z \in S$ we have $\operatorname{divisor}(f+g,D)(z) = 0$.
\end{lemma}
\begin{proof}
    \uses{lem:fg-analytic-sphere, lem:divisor-eq-zero-of-orderZero}
    Fix $z \in S \subseteq D$. By Lemma~\ref{lem:fg-analytic-sphere} the order of
    $f+g$ at $z$ is $0$, so Lemma~\ref{lem:divisor-eq-zero-of-orderZero} gives
    $\operatorname{divisor}(f+g,D)(z) = 0$.
\end{proof}

\begin{lemma}[The divisors are supported in the open disk]
    \label{lem:divisor-support-open}
    \lean{LeanEval.ComplexAnalysis.divisor_support_open}
    \leanfile{LeanEval/ComplexAnalysis/Rouche/Support.lean}
    \uses{lem:divisor-f-zero-sphere, lem:divisor-fg-zero-sphere}
    \leanok
    The divisors of $f$ and of $f+g$ on the closed disk $D$ are supported in the
    open disk $B$: both $\operatorname{supp}\operatorname{divisor}(f,D)$ and
    $\operatorname{supp}\operatorname{divisor}(f+g,D)$ are contained in $B$.
\end{lemma}
\begin{proof}
    \uses{lem:divisor-f-zero-sphere, lem:divisor-fg-zero-sphere}
    A support point $z$ of either divisor lies in $D$. If $\|z\| = R$ then
    $z \in S$, where the divisor vanishes by
    Lemmas~\ref{lem:divisor-f-zero-sphere} and~\ref{lem:divisor-fg-zero-sphere},
    contradicting $z$ being a support point. Hence $\|z\| < R$, i.e.\ $z \in B$.
\end{proof}

\subsection{The argument principle}

\begin{lemma}[Logarithmic derivative of a single power factor]
    \label{lem:logDeriv-zpow-shift}
    \lean{LeanEval.ComplexAnalysis.logDeriv_zpow_shift}
    \leanfile{LeanEval/ComplexAnalysis/Rouche/Support.lean}
    \leanok
    For $u \in \mathbb{C}$, an integer $n$, and $z \neq u$,
    \[
        \operatorname{logDeriv}\big((\cdot - u)^{n}\big)(z) = n\,(z-u)^{-1}.
    \]
\end{lemma}
\begin{proof}
    The map $z \mapsto z - u$ has derivative $1$ and is nonzero at $z$. The
    logarithmic derivative of an integer power is the exponent times the
    logarithmic derivative of the base, so
    $\operatorname{logDeriv}\big((\cdot-u)^{n}\big)(z) = n\,(z-u)^{-1}$.
\end{proof}

\begin{lemma}[Logarithmic derivative of a factorized rational function]
    \label{lem:logDeriv-factorizedRational}
    \lean{LeanEval.ComplexAnalysis.logDeriv_factorizedRational}
    \leanfile{LeanEval/ComplexAnalysis/Rouche/Support.lean}
    \uses{lem:logDeriv-zpow-shift}
    \leanok
    Let $d$ be an integer-valued divisor with finite support $T$, and put
    $\varphi = \operatorname*{{\textstyle\prod}^{\mathrm f}}_{u}(\cdot - u)^{d(u)}$,
    the finite product (\texttt{finprod}, i.e.\ the
    \texttt{Function.FactorizedRational} of $d$) over all $u$. For every
    $z \notin T$,
    \[
        \operatorname{logDeriv}\varphi(z)
          = \operatorname*{{\textstyle\sum}^{\mathrm f}}_{u} d(u)\,(z-u)^{-1},
    \]
    the finite sum (\texttt{finsum}) over the support $T$.
\end{lemma}
\begin{proof}
    \uses{lem:logDeriv-zpow-shift}
    Rewriting the \texttt{finprod} as the \texttt{Finset} product
    $\prod_{u\in T}(\cdot-u)^{d(u)}$ over the support and noting that, for
    $z \notin T$, each factor $(\cdot-u)^{d(u)}$ ($u \in T$) is differentiable and
    nonzero at $z$, \texttt{logDeriv\_prod} distributes the logarithmic derivative
    of the product as the sum over $T$ of the factors' logarithmic derivatives.
    Each summand equals $d(u)(z-u)^{-1}$ by
    Lemma~\ref{lem:logDeriv-zpow-shift}, and the resulting Finset sum is the
    \texttt{finsum} over $T$.
\end{proof}

\begin{lemma}[Each simple pole is circle-integrable]
    \label{lem:circleIntegrable-sub-inv}
    \lean{LeanEval.ComplexAnalysis.circleIntegrable_sub_inv}
    \leanfile{LeanEval/ComplexAnalysis/Rouche/Support.lean}
    \leanok
    For $u \in B$, the function $z \mapsto (z-u)^{-1}$ is circle-integrable on
    $C(0,R)$.
\end{lemma}
\begin{proof}
    The function $z \mapsto (z-u)^{-1}$ is circle-integrable on $C(0,R)$ if and
    only if $u \notin S$. Since $u \in B$ and $B \cap S = \emptyset$, indeed
    $u \notin S$.
\end{proof}

\begin{lemma}[Cauchy integral of a single simple pole]
    \label{lem:circleIntegral-sub-inv}
    \lean{LeanEval.ComplexAnalysis.circleIntegral_sub_inv}
    \leanfile{LeanEval/ComplexAnalysis/Rouche/Support.lean}
    \leanok
    For $u \in B$,
    $\oint_{C(0,R)} (z-u)^{-1}\,dz = 2\pi i$.
\end{lemma}
\begin{proof}
    Since $u \in B = B(0,R)$ we have $\|u\| < R$, so Cauchy's integral formula for
    the reciprocal on the circle $C(0,R)$ gives
    $\oint_{C(0,R)} (z-u)^{-1}\,dz = 2\pi i$.
\end{proof}

\begin{lemma}[Contour integral of a sum of simple poles]
    \label{lem:circleIntegral-sum-inv}
    \lean{LeanEval.ComplexAnalysis.circleIntegral_sum_inv}
    \leanfile{LeanEval/ComplexAnalysis/Rouche/Support.lean}
    \uses{lem:circleIntegrable-sub-inv, lem:circleIntegral-sub-inv}
    \leanok
    Let $d$ be an integer-valued divisor with finite support $T \subseteq B$.
    Then
    \[
        \oint_{C(0,R)}
          \operatorname*{{\textstyle\sum}^{\mathrm f}}_{u} d(u)\,(z-u)^{-1}\,dz
            = \Big(\operatorname*{{\textstyle\sum}^{\mathrm f}}_{u} d(u)\Big)\,
              2\pi i,
    \]
    both finite sums (\texttt{finsum}) being over the support $T$.
\end{lemma}
\begin{proof}
    \uses{lem:circleIntegrable-sub-inv, lem:circleIntegral-sub-inv}
    Each summand $z \mapsto (z-u)^{-1}$ ($u \in T \subseteq B$) is
    circle-integrable by Lemma~\ref{lem:circleIntegrable-sub-inv}, so by linearity
    the circle integral of the finite sum is the sum of the integrals and the
    constants $d(u)$ pull out. For each $u$, Lemma~\ref{lem:circleIntegral-sub-inv}
    gives $\oint_{C(0,R)} (z-u)^{-1}\,dz = 2\pi i$. Summing over $T$ yields the
    claim.
\end{proof}

\begin{lemma}[Logarithmic derivative of a unit is analytic on the disk]
  \label{lem:logDeriv-analytic-unit}
  \lean{LeanEval.ComplexAnalysis.logDeriv_analytic_unit}
  \leanfile{LeanEval/ComplexAnalysis/Rouche/Support.lean}
  \leanok
  Let U be analytic on the closed disk D and nowhere zero on D. Then
  logDeriv U = U'/U is analytic on D.
\end{lemma}
\begin{proof}
  Since U is analytic on D, so is its derivative U'. As U is analytic and
  nowhere zero on D, the quotient U'/U is analytic on D.
\end{proof}

\begin{lemma}[Contour integral of the logarithmic derivative of a unit]
    \label{lem:circleIntegral-logDeriv-unit}
    \lean{LeanEval.ComplexAnalysis.circleIntegral_logDeriv_unit}
    \leanfile{LeanEval/ComplexAnalysis/Rouche/Support.lean}
    \uses{lem:logDeriv-analytic-unit}
    Let $U$ be analytic and nowhere zero on the closed disk $D$. Then
    $\oint_{C(0,R)} \operatorname{logDeriv} U = 0$.
\end{lemma}
\begin{proof}
    By Lemma~\ref{lem:logDeriv-analytic-unit}, $\operatorname{logDeriv} U$ is
    analytic on the closed disk $D$, hence continuous on $D$ and
    complex-differentiable on the open disk (\texttt{DiffContOnCl}). By the
    Cauchy--Goursat theorem for a disk
    (\texttt{DiffContOnCl.circleIntegral\_eq\_zero}) its contour integral over
    $C(0,R)$ vanishes.
\end{proof}

\begin{lemma}[The closed disk is preperfect]
  \label{lem:closedBall-preperfect}
  \lean{LeanEval.ComplexAnalysis.closedBall_preperfect}
  \leanfile{LeanEval/ComplexAnalysis/Rouche/Support.lean}
  For R > 0 the closed disk D = closedBall(0,R) in the complex plane is
  preperfect: every point of D is an accumulation point of D.
\end{lemma}
\begin{proof}
  The open disk is open in the complex plane, which is a perfect space, so the
  open disk is preperfect (IsOpen.preperfect). Its closure is the closed disk D
  (since R > 0), and the closure of a preperfect set is perfect
  (Preperfect.perfect_closure), hence preperfect (Perfect.acc).
\end{proof}

\begin{lemma}[The two sides of the factorization are meromorphic on the disk]
  \label{lem:factorization-meromorphicAt}
  \lean{LeanEval.ComplexAnalysis.factorization_meromorphicAt}
  \leanfile{LeanEval/ComplexAnalysis/Rouche/Support.lean}
  Let h be meromorphic on the complex plane, let phi be the factorized rational
  of divisor(h,D), and let U be analytic on (a neighbourhood of) the closed disk
  D. Then for every z in D both h and phi*U are meromorphic at z.
\end{lemma}
\begin{proof}
  h is meromorphic at z since it is meromorphic on the whole plane. The
  factorized rational phi is a finite product of integer powers of the affine
  maps (.-u), each meromorphic, so phi is meromorphic at z; U is analytic at z
  (its analyticity holds on a neighbourhood of D), hence meromorphic at z; a
  product/smul of meromorphic functions is meromorphic, so phi*U is meromorphic
  at z.
\end{proof}

\begin{lemma}[The factorization holds on a full neighbourhood of the circle]
    \label{lem:factorization-eventuallyEq}
    \lean{LeanEval.ComplexAnalysis.factorization_eventuallyEq}
    \leanfile{LeanEval/ComplexAnalysis/Rouche/Support.lean}
    \uses{lem:divisor-eq-zero-of-orderZero, lem:closedBall-preperfect, lem:factorization-meromorphicAt}
    Let $h$ be meromorphic with $\operatorname{order}(h) = 0$ on $S$, let
    $\varphi = \operatorname*{{\textstyle\prod}^{\mathrm f}}_{u}
      (\cdot-u)^{\operatorname{divisor}(h,D)(u)}$ be the factorized rational
    (\texttt{finprod}) of $\operatorname{divisor}(h,D)$, and let $U$ be analytic
    and nowhere zero on $D$ with $h = \varphi \cdot U$ off a codiscrete set of
    $D$. Then for every $z \in S$ the functions $h$ and $\varphi\cdot U$ agree on
    a punctured neighbourhood of $z$.
\end{lemma}
\begin{proof}
    \uses{lem:divisor-eq-zero-of-orderZero, lem:closedBall-preperfect, lem:factorization-meromorphicAt}
    For $z \in S$ we have $z$ in the closed disk $D$. By
    Lemma~\ref{lem:factorization-meromorphicAt} both $h$ and $\varphi\cdot U$ are
    meromorphic at $z$, by Lemma~\ref{lem:closedBall-preperfect} $D$ is
    preperfect, and by hypothesis $h = \varphi U$ off a set codiscrete in $D$.
    The codiscrete-to-punctured-neighbourhood upgrade for meromorphic functions
    (\texttt{MeromorphicAt.eventuallyEq\_nhdsNE\_of\_eventuallyEq\_codiscreteWithin\_preperfect})
    then gives $h = \varphi U$ on a punctured neighbourhood of $z$.
\end{proof}

\begin{lemma}[Factorization of the logarithmic derivative on the circle]
    \label{lem:factorization-logDeriv}
    \lean{LeanEval.ComplexAnalysis.factorization_logDeriv}
    \leanfile{LeanEval/ComplexAnalysis/Rouche/Support.lean}
    \uses{lem:factorization-eventuallyEq}
    With $h$, $\varphi$, $U$ as in
    Lemma~\ref{lem:factorization-eventuallyEq}, for every $z \in S$,
    \[
        \operatorname{logDeriv} h(z)
          = \operatorname{logDeriv}\varphi(z) + \operatorname{logDeriv} U(z).
    \]
\end{lemma}
\begin{proof}
    By Lemma~\ref{lem:factorization-eventuallyEq}, $h = \varphi U$ on a punctured
    neighbourhood of $z$ (eventually with respect to $\operatorname{nhdsNE} z$).
    Both sides are analytic at $z$: $h$ has order $0$ at $z \in S$ by hypothesis,
    while $z \notin \operatorname{supp}\operatorname{divisor}(h,D)$ makes $\varphi$
    analytic at $z$ and $U$ is analytic on $D \ni z$. Two functions continuous at
    $z$ that agree on $\operatorname{nhdsNE} z$ must also agree at $z$ (their values
    there are the common limit), hence agree on the full neighbourhood
    $\operatorname{nhds} z$; therefore
    $\operatorname{logDeriv} h(z) = \operatorname{logDeriv}(\varphi U)(z)$. At $z$
    both $\varphi$ and $U$ are nonzero and differentiable, so the product rule for
    logarithmic derivatives (\texttt{logDeriv\_mul}) gives
    $\operatorname{logDeriv}(\varphi U)(z) = \operatorname{logDeriv}\varphi(z) +
    \operatorname{logDeriv} U(z)$.
\end{proof}

\begin{lemma}[A bound strictly above a finite set of larger reals]
    \label{lem:exists-lt-of-finite}
    \lean{LeanEval.ComplexAnalysis.exists_lt_of_finite}
    \leanfile{LeanEval/ComplexAnalysis/Rouche/Support.lean}
    \leanok
    Let $A$ be a finite set of real numbers with $a > R$ for every $a \in A$.
    Then there is $R' > R$ with $R' \le a$ for all $a \in A$.
\end{lemma}
\begin{proof}
    If $A$ is empty, take $R' = R + 1$. Otherwise the minimum $R' = \min A$ is
    attained, so $R' \in A$ gives $R' > R$, and $R' \le a$ for every $a \in A$ by
    minimality.
\end{proof}

\begin{lemma}[A zero/pole-free annulus outside the circle]
    \label{lem:annulus-zero-free}
    \lean{LeanEval.ComplexAnalysis.annulus_zero_free}
    \leanfile{LeanEval/ComplexAnalysis/Rouche/Support.lean}
    \uses{lem:divisor-eq-zero-of-orderZero, lem:exists-lt-of-finite}
    \leanok
    Let $h$ be meromorphic on $\mathbb{C}$ with $\operatorname{order}(h) = 0$ on
    $S$. Then there is $R' > R$ such that $\operatorname{divisor}(h, B(0,R'))$ is
    supported in the open disk $B = B(0,R)$; equivalently, $h$ has order $0$ at
    every $z$ with $R \le \|z\| < R'$.
\end{lemma}
\begin{proof}
    \uses{lem:divisor-eq-zero-of-orderZero, lem:exists-lt-of-finite}
    On the compact closed disk $\overline{B}(0,R+1)$ the divisor of $h$ has finite
    support. Any support point $z$ with $\|z\| \ge R$ in fact has $\|z\| > R$:
    were $\|z\| = R$ then $z \in S$ and the order of $h$ at $z$ is $0$, so
    Lemma~\ref{lem:divisor-eq-zero-of-orderZero} gives
    $\operatorname{divisor}(h,\cdot)(z) = 0$, contradicting $z$ being in the
    support. The moduli of these finitely many support points all exceed $R$, so
    Lemma~\ref{lem:exists-lt-of-finite} produces $R'$ with $R < R' \le R+1$ below
    all of them; then no support point lies in the annulus $R \le \|z\| < R'$.
    Restricting (\texttt{MeromorphicOn.divisor\_restrict}) shows
    $\operatorname{divisor}(h, B(0,R'))$ is supported in $B$.
\end{proof}

\begin{lemma}[A nonzero meromorphic function has finite order everywhere]
    \label{lem:order-ne-top}
    \lean{LeanEval.ComplexAnalysis.order_ne_top}
    \leanfile{LeanEval/ComplexAnalysis/Rouche/Support.lean}
    \leanok
    Let $h$ be meromorphic on $\mathbb{C}$, not identically zero. Then for every
    $z \in \mathbb{C}$ the meromorphic order of $h$ at $z$ is finite
    ($\neq \infty$). In particular, on any subset its order is finite at every
    point.
\end{lemma}
\begin{proof}
    The domain $\mathbb{C}$ is preconnected, and on a preconnected set a
    meromorphic function has order $\infty$ at some point iff it is locally
    eventually zero (equivalently, identically zero on the connected domain).
    Since $h$ is not identically zero, its order is finite at every point.
\end{proof}

\begin{lemma}[The factorized rational of $h$ on $B'$ restricts to $\varphi$]
    \label{lem:factorizedRational-restrict}
    \lean{LeanEval.ComplexAnalysis.factorizedRational_restrict}
    \leanfile{LeanEval/ComplexAnalysis/Rouche/Support.lean}
    \leanok
    Let $h$ be meromorphic on $\mathbb{C}$ and let $D \subseteq B'$ with
    $\operatorname{divisor}(h, B')$ supported in $B \subseteq D$. Then
    \[
        \operatorname*{{\textstyle\prod}^{\mathrm f}}_{u}
          (\cdot-u)^{\operatorname{divisor}(h,B')(u)}
        = \operatorname*{{\textstyle\prod}^{\mathrm f}}_{u}
          (\cdot-u)^{\operatorname{divisor}(h,D)(u)} =: \varphi .
    \]
\end{lemma}
\begin{proof}
    By \texttt{MeromorphicOn.divisor\_restrict}, restricting
    $\operatorname{divisor}(h, B')$ to $D \subseteq B'$ yields
    $\operatorname{divisor}(h, D)$. Since the support of
    $\operatorname{divisor}(h, B')$ already lies in $B \subseteq D$, restriction
    leaves it unchanged, so $\operatorname{divisor}(h, B')$ and
    $\operatorname{divisor}(h, D)$ agree as integer-valued functions. Hence their
    factorized rationals (the \texttt{finprod} over all $u$ of
    $(\cdot-u)^{\bullet(u)}$) coincide.
\end{proof}

\begin{lemma}[Existence of the factorization $h = \varphi\cdot U$]
    \label{lem:factorization-existence}
    \lean{LeanEval.ComplexAnalysis.factorization_existence}
    \leanfile{LeanEval/ComplexAnalysis/Rouche/Support.lean}
    \uses{lem:annulus-zero-free, lem:order-ne-top, lem:factorizedRational-restrict}
    \leanok
    Let $h$ be meromorphic on $\mathbb{C}$, not identically zero, with
    $\operatorname{order}(h) = 0$ on $S$. Write
    $\varphi = \operatorname*{{\textstyle\prod}^{\mathrm f}}_{u}
      (\cdot-u)^{\operatorname{divisor}(h,D)(u)}$, the factorized rational
    (\texttt{finprod}) of $\operatorname{divisor}(h,D)$. Then there is a function
    $U$, analytic and nowhere zero on the closed disk $D$, with
    $h = \varphi\cdot U$ off a codiscrete set of $D$.
\end{lemma}
\begin{proof}
    \uses{lem:annulus-zero-free, lem:order-ne-top, lem:factorizedRational-restrict}
    Choose $R' > R$ as in Lemma~\ref{lem:annulus-zero-free}, so the open disk
    $B' = B(0,R')$ contains $D$ and $\operatorname{divisor}(h, B')$ is supported
    in $B$. By Lemma~\ref{lem:order-ne-top} the order of $h$ is finite at every
    point of $B'$, so the extraction of zeros and poles on $B'$
    (\texttt{MeromorphicOn.extract\_zeros\_poles}) produces an analytic
    nowhere-zero function $g$ on $B'$ with
    $h = \big(\operatorname*{{\textstyle\prod}^{\mathrm f}}_{u}
    (\cdot-u)^{\operatorname{divisor}(h,B')(u)}\big)\cdot g$ off a set codiscrete
    in $B'$. By Lemma~\ref{lem:factorizedRational-restrict} that factorized
    rational equals $\varphi$. Take $U = g$ restricted to $D$: it is analytic and
    nowhere zero on $D \subseteq B'$ (\texttt{AnalyticOnNhd.mono}), and the
    codiscrete equality on $B'$ restricts to one on $D$
    (\texttt{Filter.codiscreteWithin\_mono}).
\end{proof}

\begin{lemma}[$\operatorname{logDeriv}\varphi$ is circle-integrable]
    \label{lem:circleIntegrable-logDeriv-factorizedRational}
    \lean{LeanEval.ComplexAnalysis.circleIntegrable_logDeriv_factorizedRational}
    \leanfile{LeanEval/ComplexAnalysis/Rouche/Support.lean}
    \uses{lem:logDeriv-factorizedRational, lem:circleIntegrable-sub-inv}
    Let $\varphi = \operatorname*{{\textstyle\prod}^{\mathrm f}}_{u}
      (\cdot-u)^{\operatorname{divisor}(h,D)(u)}$ be the factorized rational of
    $\operatorname{divisor}(h,D)$, whose finite support $T$ lies in $B$. Then
    $\operatorname{logDeriv}\varphi$ is circle-integrable on $C(0,R)$.
\end{lemma}
\begin{proof}
    \uses{lem:logDeriv-factorizedRational, lem:circleIntegrable-sub-inv}
    For $z \in S$ we have $z \notin T$ (as $T \subseteq B$), so by
    Lemma~\ref{lem:logDeriv-factorizedRational} the function
    $\operatorname{logDeriv}\varphi$ agrees on the circle with the finite sum
    $\sum_{u}^{\mathrm f}\operatorname{divisor}(h,D)(u)(z-u)^{-1}$. Each summand is
    a constant multiple of $z \mapsto (z-u)^{-1}$, circle-integrable by
    Lemma~\ref{lem:circleIntegrable-sub-inv} (every $u \in T \subseteq B$), so the
    finite sum is circle-integrable, and circle-integrability transfers across the
    pointwise equality on the circle (\texttt{circleIntegrable\_congr}).
\end{proof}

\begin{lemma}[$\operatorname{logDeriv} U$ is circle-integrable]
    \label{lem:circleIntegrable-logDeriv-unit}
    \lean{LeanEval.ComplexAnalysis.circleIntegrable_logDeriv_unit}
    \leanfile{LeanEval/ComplexAnalysis/Rouche/Support.lean}
    \uses{lem:logDeriv-analytic-unit}
    Let $U$ be analytic and nowhere zero on the closed disk $D$. Then
    $\operatorname{logDeriv} U$ is circle-integrable on $C(0,R)$.
\end{lemma}
\begin{proof}
    \uses{lem:logDeriv-analytic-unit}
    By Lemma~\ref{lem:logDeriv-analytic-unit} $\operatorname{logDeriv} U$ is
    analytic on $D$, hence continuous on the circle $S \subseteq D$. A continuous
    function on the circle is circle-integrable
    (\texttt{ContinuousOn.circleIntegrable}).
\end{proof}

\begin{lemma}[Splitting the contour integral of $\operatorname{logDeriv} h$]
    \label{lem:argument-principle-integral-split}
    \lean{LeanEval.ComplexAnalysis.argument_principle_integral_split}
    \leanfile{LeanEval/ComplexAnalysis/Rouche/Support.lean}
    \uses{lem:factorization-existence, lem:factorization-logDeriv, lem:circleIntegrable-logDeriv-factorizedRational, lem:circleIntegrable-logDeriv-unit}
    Let $h$ be meromorphic on $\mathbb{C}$, not identically zero, with
    $\operatorname{order}(h) = 0$ on $S$, and let $h = \varphi\cdot U$ be the
    factorization of Lemma~\ref{lem:factorization-existence}, with $\varphi$ the
    factorized rational of $\operatorname{divisor}(h,D)$ and $U$ analytic and
    nowhere zero on $D$. Then
    \[
        \oint_{C(0,R)}\operatorname{logDeriv} h
          = \oint_{C(0,R)}\operatorname{logDeriv}\varphi
          + \oint_{C(0,R)}\operatorname{logDeriv} U .
    \]
\end{lemma}
\begin{proof}
    \uses{lem:factorization-logDeriv, lem:circleIntegrable-logDeriv-factorizedRational, lem:circleIntegrable-logDeriv-unit}
    On $S$, Lemma~\ref{lem:factorization-logDeriv} gives the pointwise identity
    $\operatorname{logDeriv} h = \operatorname{logDeriv}\varphi +
    \operatorname{logDeriv} U$. Both summands are circle-integrable, by
    Lemmas~\ref{lem:circleIntegrable-logDeriv-factorizedRational}
    and~\ref{lem:circleIntegrable-logDeriv-unit}, so by linearity of the circle
    integral (\texttt{integral\_add}) the integral of the sum splits as claimed.
\end{proof}

\begin{lemma}[Contour integral of $\operatorname{logDeriv}\varphi$]
    \label{lem:circleIntegral-logDeriv-factorizedRational}
    \lean{LeanEval.ComplexAnalysis.circleIntegral_logDeriv_factorizedRational}
    \leanfile{LeanEval/ComplexAnalysis/Rouche/Support.lean}
    \uses{lem:logDeriv-factorizedRational, lem:circleIntegral-sum-inv}
    \leanok
    Let $\varphi = \operatorname*{{\textstyle\prod}^{\mathrm f}}_{u}
      (\cdot-u)^{\operatorname{divisor}(h,D)(u)}$ be the factorized rational of
    $\operatorname{divisor}(h,D)$, whose finite support $T$ lies in $B$. Then
    \[
        \oint_{C(0,R)}\operatorname{logDeriv}\varphi
          = \Big(\operatorname*{{\textstyle\sum}^{\mathrm f}}_{u}
            \operatorname{divisor}(h,D)(u)\Big)\, 2\pi i .
    \]
\end{lemma}
\begin{proof}
    \uses{lem:logDeriv-factorizedRational, lem:circleIntegral-sum-inv}
    For $z \in S$ we have $z \notin T$ (as $T \subseteq B$), so
    Lemma~\ref{lem:logDeriv-factorizedRational} gives
    $\operatorname{logDeriv}\varphi(z) =
    \operatorname*{{\textstyle\sum}^{\mathrm f}}_{u}
    \operatorname{divisor}(h,D)(u)(z-u)^{-1}$. Integrating over $C(0,R)$ and
    applying Lemma~\ref{lem:circleIntegral-sum-inv} yields the claim.
\end{proof}

\begin{lemma}[Integral evaluation of the logarithmic derivative]
    \label{lem:argument-principle-integral}
    \lean{LeanEval.ComplexAnalysis.argument_principle_integral}
    \leanfile{LeanEval/ComplexAnalysis/Rouche/Support.lean}
    \uses{lem:factorization-existence, lem:argument-principle-integral-split, lem:circleIntegral-logDeriv-factorizedRational, lem:circleIntegral-logDeriv-unit}
    \leanok
    Let $h$ be meromorphic on $\mathbb{C}$, not identically zero, with
    $\operatorname{order}(h) = 0$ on $S$. Then
    \[
        \oint_{C(0,R)}\operatorname{logDeriv} h
          = \Big(\operatorname*{{\textstyle\sum}^{\mathrm f}}_{u}
            \operatorname{divisor}(h,D)(u)\Big)\, 2\pi i .
    \]
\end{lemma}
\begin{proof}
    \uses{lem:factorization-existence, lem:argument-principle-integral-split,
      lem:circleIntegral-logDeriv-factorizedRational,
      lem:circleIntegral-logDeriv-unit}
    By Lemma~\ref{lem:factorization-existence} write $h = \varphi\cdot U$ with
    $\varphi$ the factorized rational of $\operatorname{divisor}(h,D)$ and $U$
    analytic and nowhere zero on $D$. By
    Lemma~\ref{lem:argument-principle-integral-split},
    $\oint_{C(0,R)}\operatorname{logDeriv} h =
    \oint_{C(0,R)}\operatorname{logDeriv}\varphi +
    \oint_{C(0,R)}\operatorname{logDeriv} U$. The second integral vanishes by
    Lemma~\ref{lem:circleIntegral-logDeriv-unit}, and the first equals
    $\big(\operatorname*{{\textstyle\sum}^{\mathrm f}}_{u}
    \operatorname{divisor}(h,D)(u)\big)2\pi i$ by
    Lemma~\ref{lem:circleIntegral-logDeriv-factorizedRational}.
\end{proof}

\begin{theorem}[Argument principle, zero-counting form]
    \label{thm:argument-principle}
    \lean{LeanEval.ComplexAnalysis.argument_principle}
    \leanfile{LeanEval/ComplexAnalysis/Rouche/Support.lean}
    \uses{lem:argument-principle-integral}
    \leanok
    Let $h$ be meromorphic on $\mathbb{C}$, not identically zero, and assume that
    on $S$ the order of $h$ is $0$ (so $h$ is analytic and nonzero on the
    circle). Then
    \[
        \sum_{z}\operatorname{divisor}(h,D)(z)
          = \frac{1}{2\pi i}\oint_{C(0,R)} \operatorname{logDeriv} h .
    \]
\end{theorem}
\begin{proof}
    The total mass $\sum_z \operatorname{divisor}(h,D)(z)$ is by definition the
    finite sum (\texttt{finsum})
    $\operatorname*{{\textstyle\sum}^{\mathrm f}}_{u}\operatorname{divisor}(h,D)(u)$
    over the support. By Lemma~\ref{lem:argument-principle-integral},
    $\oint_{C(0,R)}\operatorname{logDeriv} h =
    \big(\operatorname*{{\textstyle\sum}^{\mathrm f}}_{u}
    \operatorname{divisor}(h,D)(u)\big)2\pi i$. Dividing both sides by $2\pi i$
    gives the claim.
\end{proof}

\subsection{Vanishing winding of \texorpdfstring{$(f+g)/f$}{(f+g)/f}}

\begin{lemma}[The quotient maps the circle into a slit-free disk]
    \label{lem:F-slitPlane}
    \lean{LeanEval.ComplexAnalysis.F_slitPlane}
    \leanfile{LeanEval/ComplexAnalysis/Rouche/Support.lean}
    \uses{lem:f-ne-zero-sphere}
    \leanok
    For every $z \in S$, the value $F(z) := (f+g)(z)/f(z)$ satisfies
    $\|F(z) - 1\| < 1$; in particular $F(z)$ lies in the slit plane
    $\mathbb{C}\setminus(-\infty,0]$.
\end{lemma}
\begin{proof}
    By Lemma~\ref{lem:f-ne-zero-sphere} we have $f(z) \neq 0$, so $F(z) - 1 =
    g(z)/f(z)$ and $\|F(z)-1\| = \|g(z)\|/\|f(z)\| < 1$ by the hypothesis. Thus
    $F(z) \in B(1,1)$, and the open ball $B(1,1)$ is contained in the slit plane
    $\mathbb{C}\setminus(-\infty,0]$.
\end{proof}

\begin{lemma}[$F = (f+g)/f$ is differentiable and nonzero on the circle]
    \label{lem:F-hasDerivAt}
    \lean{LeanEval.ComplexAnalysis.F_hasDerivAt}
    \leanfile{LeanEval/ComplexAnalysis/Rouche/Support.lean}
    \uses{lem:f-analytic-sphere, lem:fg-analytic-sphere}
    \leanok
    With $F = (f+g)/f$, for every $z \in S$ the function $F$ is differentiable at
    $z$ — so $\operatorname{HasDerivAt} F\,(F'(z))\,z$ for its derivative
    $F'(z)$ — and $F(z) \neq 0$.
\end{lemma}
\begin{proof}
    By Lemmas~\ref{lem:f-analytic-sphere} and~\ref{lem:fg-analytic-sphere}, near
    $z$ both $f$ and $f+g$ are analytic with nonzero values. Hence the quotient
    $F = (f+g)/f$ is differentiable at $z$, with derivative $F'(z)$ given by the
    quotient rule, and $F(z) = (f+g)(z)/f(z) \neq 0$.
\end{proof}

\begin{lemma}[The logarithm of $F$ is a primitive of $\operatorname{logDeriv} F$ on the circle]
    \label{lem:winding-hasDerivWithinAt}
    \lean{LeanEval.ComplexAnalysis.winding_hasDerivWithinAt}
    \leanfile{LeanEval/ComplexAnalysis/Rouche/Support.lean}
    \uses{lem:F-slitPlane, lem:F-hasDerivAt}
    \leanok
    With $F = (f+g)/f$, for every $z \in S$ the function
    $\operatorname{Log}\circ F$ has derivative $\operatorname{logDeriv} F(z)$
    within $S$ at $z$.
\end{lemma}
\begin{proof}
    \uses{lem:F-slitPlane, lem:F-hasDerivAt}
    By Lemma~\ref{lem:F-hasDerivAt}, $\operatorname{HasDerivAt} F\,(F'(z))\,z$ and
    $F(z) \neq 0$. By Lemma~\ref{lem:F-slitPlane}, $F(z)$ lies in the slit plane,
    where $\operatorname{Log}$ is differentiable with derivative $F(z)^{-1}$. The
    chain rule for $\operatorname{Log}$ (\texttt{HasDerivAt.clog}) then gives
    $\operatorname{HasDerivAt}(\operatorname{Log}\circ F)(F'(z)/F(z))\,z$; since
    $F'(z)/F(z) = \operatorname{logDeriv} F(z)$, restricting to $S$ yields the
    within-$S$ derivative.
\end{proof}

\begin{lemma}[The winding contour integral vanishes]
    \label{lem:winding-cancel}
    \lean{LeanEval.ComplexAnalysis.winding_cancel}
    \leanfile{LeanEval/ComplexAnalysis/Rouche/Support.lean}
    \uses{lem:winding-hasDerivWithinAt}
    \leanok
    With $F = (f+g)/f$, $\oint_{C(0,R)} \operatorname{logDeriv} F = 0$.
\end{lemma}
\begin{proof}
    By Lemma~\ref{lem:winding-hasDerivWithinAt}, at every $z \in S$ the value
    $\operatorname{logDeriv} F(z)$ is the derivative within $S$ of
    $\operatorname{Log}\circ F$. As $R \ge 0$, the circle integral of a function
    that is, everywhere on the circle, such a within-circle derivative vanishes.
\end{proof}

\begin{lemma}[Pointwise split of $\operatorname{logDeriv}((f+g)/f)$ on the circle]
    \label{lem:logDeriv-div-pointwise}
    \lean{LeanEval.ComplexAnalysis.logDeriv_div_pointwise}
    \leanfile{LeanEval/ComplexAnalysis/Rouche/Support.lean}
    \uses{lem:f-analytic-sphere, lem:fg-analytic-sphere}
    \leanok
    For every $z \in S$,
    $\operatorname{logDeriv}\big((f+g)/f\big)(z)
       = \operatorname{logDeriv}(f+g)(z) - \operatorname{logDeriv} f(z)$.
\end{lemma}
\begin{proof}
    On $S$ both $f$ and $f+g$ are analytic and nonzero
    (Lemmas~\ref{lem:f-analytic-sphere} and~\ref{lem:fg-analytic-sphere}); in
    particular both are differentiable at $z$ with nonzero values. Hence
    \texttt{logDeriv\_div} applies and the logarithmic derivative of the quotient
    is the difference of the logarithmic derivatives.
\end{proof}

\begin{lemma}[The logarithmic derivatives are circle-integrable]
    \label{lem:logDeriv-circleIntegrable}
    \lean{LeanEval.ComplexAnalysis.logDeriv_circleIntegrable}
    \leanfile{LeanEval/ComplexAnalysis/Rouche/Support.lean}
    \uses{lem:f-analytic-sphere, lem:fg-analytic-sphere}
    \leanok
    The functions $\operatorname{logDeriv} f$ and $\operatorname{logDeriv}(f+g)$
    are circle-integrable on $C(0,R)$.
\end{lemma}
\begin{proof}
    \uses{lem:f-analytic-sphere, lem:fg-analytic-sphere}
    By Lemmas~\ref{lem:f-analytic-sphere} and~\ref{lem:fg-analytic-sphere}, each of
    $f$ and $f+g$ is analytic and nonzero at every $z \in S$. Hence
    $\operatorname{logDeriv} f = f'/f$ and $\operatorname{logDeriv}(f+g)$ are
    continuous on the circle, and a continuous function on the circle is
    circle-integrable (\texttt{ContinuousOn.circleIntegrable}).
\end{proof}

\begin{lemma}[Equality of the two winding integrals]
    \label{lem:logDeriv-diff}
    \lean{LeanEval.ComplexAnalysis.logDeriv_diff}
    \leanfile{LeanEval/ComplexAnalysis/Rouche/Support.lean}
    \uses{lem:winding-cancel, lem:logDeriv-div-pointwise, lem:logDeriv-circleIntegrable}
    \leanok
    $\displaystyle \oint_{C(0,R)} \operatorname{logDeriv}(f+g)
       = \oint_{C(0,R)} \operatorname{logDeriv} f$.
\end{lemma}
\begin{proof}
    \uses{lem:winding-cancel, lem:logDeriv-div-pointwise,
      lem:logDeriv-circleIntegrable}
    By Lemma~\ref{lem:logDeriv-circleIntegrable} both
    $\operatorname{logDeriv}(f+g)$ and $\operatorname{logDeriv} f$ are
    circle-integrable, so by \texttt{integral\_sub}
    \[
        \oint_{C(0,R)}\operatorname{logDeriv}(f+g)
          - \oint_{C(0,R)}\operatorname{logDeriv} f
          = \oint_{C(0,R)}\big(\operatorname{logDeriv}(f+g)
            - \operatorname{logDeriv} f\big).
    \]
    By Lemma~\ref{lem:logDeriv-div-pointwise} the integrand agrees on $S$ with
    $\operatorname{logDeriv}\big((f+g)/f\big)$, whose circle integral is $0$ by
    Lemma~\ref{lem:winding-cancel}. Hence the difference of the two integrals is
    $0$.
\end{proof}

\subsection{Matching pole divisors and conclusion}

\begin{lemma}[Nonnegative order is preserved by adding $g$]
    \label{lem:order-nonneg-of-pole}
    \lean{LeanEval.ComplexAnalysis.order_nonneg_of_pole}
    \leanfile{LeanEval/ComplexAnalysis/Rouche/Support.lean}
    \leanok
    For every $z \in \mathbb{C}$, if $\operatorname{order}(f) \ge 0$ at $z$ then
    $\operatorname{order}(f+g) \ge 0$ at $z$.
\end{lemma}
\begin{proof}
    Since $g$ is analytic, $\operatorname{order}(g) \ge 0$ at $z$. If
    $\operatorname{order}(f) \ge 0$, then $\operatorname{order}(f+g) \ge
    \min(\operatorname{order} f, \operatorname{order} g) \ge 0$.
\end{proof}

\begin{lemma}[Order of $f+g$ equals that of $f$ at a pole of $f$]
    \label{lem:order-eq-at-pole}
    \lean{LeanEval.ComplexAnalysis.order_eq_at_pole}
    \leanfile{LeanEval/ComplexAnalysis/Rouche/Support.lean}
    \leanok
    For every $z \in \mathbb{C}$, if $\operatorname{order}(f) < 0$ at $z$ then
    $\operatorname{order}(f+g) = \operatorname{order}(f)$ at $z$.
\end{lemma}
\begin{proof}
    Since $g$ is analytic, $\operatorname{order}(g) \ge 0$ at $z$. If
    $\operatorname{order}(f) < 0 \le \operatorname{order}(g)$, then
    $\operatorname{order}(f) < \operatorname{order}(g)$, so the order of the sum
    equals the strictly smaller one:
    $\operatorname{order}(f+g) = \operatorname{order}(f)$.
\end{proof}

\begin{lemma}[$f$ and $f+g$ have the same pole divisor]
    \label{lem:negPart-eq}
    \lean{LeanEval.ComplexAnalysis.negPart_eq}
    \leanfile{LeanEval/ComplexAnalysis/Rouche/Support.lean}
    \uses{lem:order-nonneg-of-pole, lem:order-eq-at-pole}
    \leanok
    For every $z \in D$,
    $\big(\operatorname{divisor}(f+g,D)\big)^{-}(z)
       = \big(\operatorname{divisor}(f,D)\big)^{-}(z)$.
\end{lemma}
\begin{proof}
    \uses{lem:order-nonneg-of-pole, lem:order-eq-at-pole}
    Fix $z \in D$. The divisor at $z$ is the untopped meromorphic order, and the
    negative part of an integer is determined by that integer. If
    $\operatorname{order}(f) \ge 0$ at $z$, then by
    Lemma~\ref{lem:order-nonneg-of-pole} also $\operatorname{order}(f+g) \ge 0$,
    so both divisors are $\ge 0$ and both negative parts vanish. If
    $\operatorname{order}(f) < 0$, then by Lemma~\ref{lem:order-eq-at-pole} the
    two orders are equal, hence so are the two divisors and their negative parts.
\end{proof}

\begin{lemma}[The argument principle applies to $f$ and to $f+g$]
    \label{lem:orderZero-bundle}
    \lean{LeanEval.ComplexAnalysis.orderZero_bundle}
    \leanfile{LeanEval/ComplexAnalysis/Rouche/Support.lean}
    \uses{lem:f-ne-zero-sphere, lem:fg-ne-zero-sphere, lem:f-analytic-sphere, lem:fg-analytic-sphere}
    Each of $f$ and $f+g$ is meromorphic on $\mathbb{C}$, not identically zero,
    and has order $0$ at every $z \in S$. Hence
    Theorem~\ref{thm:argument-principle} applies to each.
\end{lemma}
\begin{proof}
    $f$ is meromorphic by hypothesis and $f+g$ is meromorphic, being the sum of a
    meromorphic and an analytic function. Neither is identically zero: $f$ is
    nonzero on $S$ (Lemma~\ref{lem:f-ne-zero-sphere}) and $f+g$ is nonzero on $S$
    (Lemma~\ref{lem:fg-ne-zero-sphere}). The order is $0$ on $S$ by
    Lemmas~\ref{lem:f-analytic-sphere} and~\ref{lem:fg-analytic-sphere}. These are
    exactly the hypotheses of Theorem~\ref{thm:argument-principle}.
\end{proof}

\begin{lemma}[The positive part decomposes the total mass]
    \label{lem:posPart-mass-decomp}
    \lean{LeanEval.ComplexAnalysis.posPart_mass_decomp}
    \leanfile{LeanEval/ComplexAnalysis/Rouche/Support.lean}
    \leanok
    For a meromorphic $h$,
    \[
        \sum_z \big(\operatorname{divisor}(h,D)\big)^{+}(z)
          = \sum_z \operatorname{divisor}(h,D)(z)
          + \sum_z \big(\operatorname{divisor}(h,D)\big)^{-}(z).
    \]
\end{lemma}
\begin{proof}
    For an integer $v$ one has $v^{+} = v + v^{-}$ (i.e.\
    $v^{+} - v^{-} = v$). Applying this pointwise to
    $v = \operatorname{divisor}(h,D)(z)$ and summing over the finite support gives
    the identity.
\end{proof}

\begin{theorem}[Rouché's theorem, zero-counting form]
    \label{thm:rouche}
    \lean{LeanEval.ComplexAnalysis.rouche_zero_count_eq}
    \leanfile{LeanEval/ComplexAnalysis/Rouche.lean}
    \uses{thm:argument-principle, lem:logDeriv-diff, lem:negPart-eq, lem:orderZero-bundle, lem:posPart-mass-decomp}
    \leanok
    Under the standing hypotheses ($R>0$, $f$ meromorphic in normal form on
    $\mathbb{C}$, $g$ analytic on $\mathbb{C}$, and $\|g(z)\| < \|f(z)\|$ on
    $S$),
    \[
        \sum_{z}\big(\operatorname{divisor}(f+g,D)\big)^{+}(z)
          = \sum_{z}\big(\operatorname{divisor}(f,D)\big)^{+}(z).
    \]
\end{theorem}
\begin{proof}
    \uses{thm:argument-principle, lem:logDeriv-diff, lem:negPart-eq,
      lem:orderZero-bundle, lem:posPart-mass-decomp}
    By Lemma~\ref{lem:posPart-mass-decomp}, applied to each of $f$ and $f+g$, the
    positive-part mass decomposes as the signed mass plus the negative-part mass.
    By Lemma~\ref{lem:orderZero-bundle}, Theorem~\ref{thm:argument-principle}
    applies to both $f$ and $f+g$:
    \[
        \sum_z \operatorname{divisor}(f+g,D)(z)
          = \frac{1}{2\pi i}\oint_{C(0,R)}\operatorname{logDeriv}(f+g),
        \qquad
        \sum_z \operatorname{divisor}(f,D)(z)
          = \frac{1}{2\pi i}\oint_{C(0,R)}\operatorname{logDeriv} f.
    \]
    By Lemma~\ref{lem:logDeriv-diff} the two integrals are equal, hence the two
    signed masses are equal. By Lemma~\ref{lem:negPart-eq} the negative-part
    masses are equal as well. Adding the equal signed and negative-part masses
    gives equality of the positive-part masses, which is the claim.
\end{proof}

\end{document}
